\section{\label{apx:additional}Additional results and analyses}
\subsection{\label{apx:sensitivity}Sensitivity to training configuration}
Three configuration blocks were run on the entity-rich dataset,
nine adapters each,
to test whether the flatness of private optimization and the choice of three epochs are properties of the mechanism or of our settings (Table~\ref{tab:sensitivity}).

\begin{table}[t]
  \centering
  \caption{Configuration blocks, entity-rich dataset,
  nine adapters each.
  The grid configuration is clipping bound $1.0$,
  rank $16$,
  three epochs.}
  \label{tab:sensitivity}
  \begin{tabular}{llccc}
    \toprule
    Block & Setting & Fidelity & Instruction & $Y$ \\
    \midrule
    \multirow{3}{*}{Clipping bound}
      & $0.1$  & 0.9303 & 0.2530 & $-0.417$ \\
      & $0.3$  & 0.9300 & 0.2559 & $-0.412$ \\
      & $3.0$  & 0.9260 & 0.2481 & $-0.403$ \\
    \midrule
    \multirow{3}{*}{Adapter rank}
      & $4$    & 0.9417 & 0.2964 & $-0.596$ \\
      & $8$    & 0.9414 & 0.2777 & $-0.547$ \\
      & $32$   & 0.9022 & 0.2138 & $-0.508$ \\
    \midrule
    \multirow{2}{*}{Six epochs}
      & Baseline      & 0.9695 & 0.2627 & -- \\
      & Private optim.& 0.9007 & 0.1717 & $-0.076$ \\
    \bottomrule
  \end{tabular}
\end{table}

\paragraph{The clip is inert.}
Varying the clipping bound from $0.1$ to $3.0$ -- a thirtyfold range spanning both sides of the mean gradient norm -- moves fidelity by $0.004$ and $Y$ by $0.014$.
Together with the budget flatness of \apx{clipping} this says the mechanism's two tuning knobs are both inert at this scale,
which strengthens rather than weakens the clipping-bias account:
when most examples are clipped at every setting ($0.561$--$0.753$ of them),
moving the threshold changes which examples are clipped rather than whether clipping occurs.

\paragraph{Rank is not.}
Adapter rank moves outcomes more than either privacy knob:
instruction following falls from $0.296$ at rank $4$ to $0.214$ at rank $32$,
and $Y$ rises by $0.088$ over the same range.
A larger adapter fits the perturbed data better and is damaged more by it.
Any comparison of private-optimization cost across papers is therefore confounded by adapter capacity unless it is held fixed.

\paragraph{Three epochs against six.}
On clean text six epochs reproduces three to within $0.005$ on fidelity and $0.001$ on instruction following,
so the baseline is saturated.
Under private optimization it is not:
six epochs costs $0.09$ of instruction following relative to the three-epoch runs and raises $Y$ by $0.31$.
Longer private training costs more here,
consistent with accumulated clipping bias rather than with undertraining.

\subsection{Secondary truthfulness measures}
The suite records four truthfulness quantities beyond the primary multiple-choice accuracy (Table~\ref{tab:truthfulqa}).
All are computed on the same $817$ questions.

\begin{table}[t]
  \centering
  \caption{TruthfulQA measures by mechanism,
  averaged over datasets,
  levels,
  models and seeds.
  Normalized variants divide the option log-probability by option length;
  expected calibration error (ECE)~\citep{guo2017calibration} uses ten equal-width confidence bins.}
  \label{tab:truthfulqa}
  \small
  \setlength{\tabcolsep}{4pt}
  \begin{tabular}{lcccccc}
    \toprule
    Mechanism & MC1 & MC1$_n$ & MC2 & MC2$_n$ & Margin & ECE \\
    \midrule
    Base                 & 0.3137 & 0.3639 & 0.4904 & 0.4578 & $-0.9977$ & 0.1672 \\
    Baseline adaptation  & 0.2382 & 0.2782 & 0.3828 & 0.4017 & $-0.9025$ & 0.1636 \\
    Random replacement   & 0.2269 & 0.2711 & 0.3869 & 0.4034 & $-0.7714$ & 0.1428 \\
    Sanitization         & 0.2500 & 0.2862 & 0.4241 & 0.4189 & $-0.6678$ & 0.1343 \\
    Entity-targeted      & 0.2390 & 0.2780 & 0.3833 & 0.4056 & $-0.8177$ & 0.1535 \\
    Private optimization & 0.2331 & 0.2674 & 0.3812 & 0.4018 & $-0.8568$ & 0.1585 \\
    \bottomrule
  \end{tabular}
\end{table}

Three observations.
Length normalization raises MC1 by four to five points uniformly and does not change any ordering.
Adaptation costs about eight points of MC1 and eleven of MC2 relative to the base model,
consistent with the instruction-following loss of \apx{sanity}.
And calibration error \textit{falls} under the text-perturbing mechanisms,
from $0.164$ at the non-private baseline to $0.134$ under sanitization,
in the same direction as the rising margin:
distorted models are less confident and,
on this measure,
better calibrated.
The entity-targeted control sits at $0.154$,
between the baseline and the uniform mechanism,
tracking its intermediate carrier damage.
Whether this is a benefit or an artifact of flattened output distributions is not resolved here.

\subsection{Instruction-level accuracy}
IFEval admits a stricter prompt-level score,
requiring every instruction in a prompt to be satisfied,
and a looser instruction-level score.
Both are recorded;
the primary analysis uses the former.
They move together (Table~\ref{tab:ifevallevels}),
with the instruction-level score running $10$--$11$ points higher throughout and preserving every mechanism ordering,
so no conclusion depends on the choice.

\begin{table}[t]
  \centering
  \caption{IFEval at both granularities,
  averaged over levels,
  models and seeds.}
  \label{tab:ifevallevels}
  \begin{tabular}{llcc}
    \toprule
    Dataset & Mechanism & Prompt-level & Instruction-level \\
    \midrule
    \multirow{3}{*}{Poor}
      & Baseline       & 0.3383 & 0.4543 \\
      & Replacement    & 0.2126 & 0.3218 \\
      & Sanitization   & 0.2249 & 0.3308 \\
      & Private optim. & 0.3178 & 0.4338 \\
    \midrule
    \multirow{5}{*}{Rich}
      & Baseline       & 0.2639 & 0.3750 \\
      & Replacement    & 0.1543 & 0.2513 \\
      & Sanitization   & 0.1993 & 0.3067 \\
      & Entity-targeted& 0.2610 & 0.3713 \\
      & Private optim. & 0.2560 & 0.3688 \\
    \bottomrule
  \end{tabular}
\end{table}

\subsection{Accounting and compute}
The privacy accountant's realized budgets match their targets to three decimals ($0.991$, $3.995$, $7.999$, $15.995$ against nominal $1$, $4$, $8$, $16$),
at noise multipliers of $1.299$,
$0.740$,
$0.581$ and $0.453$ and a mean sampled lot size of $238$.
The delta is fixed below the reciprocal of the dataset size throughout.

The full grid cost approximately $407$ GPU-hours:
$339$ in training -- $33.6$ minutes for a mean non-private cell against $71.0$ for a private one,
the private branch being roughly twice as expensive per cell because per-example gradients are computed by loop rather than vectorised (\apx{dpsgd}) -- $66$ across both evaluation suites,
and under three hours for the single distortion pass over $26$ text variants and $387$ generation sets.

\subsection{Coverage}
The grid completed in full.
The entity-poor dataset contributes $17$ conditions per model -- one non-private baseline,
four levels of random replacement,
eight of sanitization and four of private optimization -- and the entity-rich dataset adds two levels of entity-targeted sanitization for $19$.
Across three base models and three seeds this is $153$ and $171$ adapters,
$324$ in total,
from $444$ training runs;
the excess is the tuning wave of \apx{calib} together with the configuration-sensitivity runs of \apx{sensitivity}.
Every adapter was evaluated on the alignment suite and,
separately,
on the refusal suite,
giving $408$ evaluations of each kind ($405$ adapters plus the three untuned base models).

Semantic distortion was scored on all $26$ perturbed text variants and on $387$ generation sets -- one for every adapter that has a non-private counterpart to be compared against,
the remainder being the $18$ non-private baselines themselves and the three untuned models.

\subsection{\label{apx:sanity}Sanity checks}
\paragraph{Null perturbation.}
Sanitization at $\varepsilon = 12$ leaves the text bit-for-bit unchanged ($S \approx 10^{-7}$),
so it must reproduce the non-private baseline.
On the entity-rich dataset it does,
to within $0.012$ on every metric and every model;
the largest discrepancy is $+0.0117$ on instruction following.
On the entity-poor dataset the gap is larger and negative -- entity fidelity falls by $0.018$,
$0.044$ and $0.137$ for the three models -- but the seed spread of that probe on that dataset is $0.13$--$0.23$ (Table~\ref{tab:seednoise}),
so the deviation sits inside the noise.
The check passes where the instrument is precise and is uninformative where it is not.

\paragraph{Seed noise.}
Table~\ref{tab:seednoise} reports the mean within-condition range across the three seeds.
Instruction following and TruthfulQA are stable on both datasets ($0.012$--$0.046$).
Entity fidelity is stable on the entity-rich dataset ($0.009$--$0.067$) and noisier on the entity-poor one ($0.132$--$0.234$):
the extraction task is held out from their training,
so the probe measures generalization rather than fidelity.
This asymmetry recurs throughout and is the single largest threat to the moderation analysis (\apx{moderation}).

\begin{table}[t]
  \centering
  \caption{Mean within-condition range across three seeds.}
  \label{tab:seednoise}
  \begin{tabular}{llccc}
    \toprule
    Dataset & Mechanism & Instruction & Fidelity & MC1 \\
    \midrule
    \multirow{4}{*}{Entity-poor}
     & Baseline             & 0.0246 & 0.1748 & 0.0208 \\
     & Random replacement   & 0.0294 & 0.1784 & 0.0136 \\
     & Sanitization         & 0.0238 & 0.1318 & 0.0135 \\
     & Private optimization & 0.0368 & 0.2344 & 0.0126 \\
    \midrule
    \multirow{5}{*}{Entity-rich}
     & Baseline             & 0.0351 & 0.0100 & 0.0188 \\
     & Random replacement   & 0.0331 & 0.0340 & 0.0186 \\
     & Sanitization         & 0.0280 & 0.0666 & 0.0221 \\
     & Entity-targeted      & 0.0354 & 0.0086 & 0.0184 \\
     & Private optimization & 0.0456 & 0.0611 & 0.0251 \\
    \bottomrule
  \end{tabular}
\end{table}

\paragraph{The cost of adaptation itself.}
Before any privacy mechanism acts,
non-private adaptation already moves the models a long way (Table~\ref{tab:basevsm0}).
Instruction following falls by $0.19$--$0.32$;
entity fidelity rises from near zero to $0.76$--$0.83$,
since the base models do not perform the extraction task at all.
Because $Y$ is defined against the non-private baseline of the same model and dataset,
this shift is absorbed and does not enter the regression -- but it sets the scale against which every mechanism effect below should be read.

\begin{table}[t]
  \centering
  \caption{Base models against their non-private adapted counterparts,
  averaged over datasets and seeds.}
  \label{tab:basevsm0}
  \begin{tabular}{llcccc}
    \toprule
    Model & State & Instruction & Fidelity & MC1 & Leakage \\
    \midrule
    \multirow{2}{*}{Falcon3-1B-Instruct}
      & Base     & 0.4824 & 0.2038 & 0.3084 & 0.0069 \\
      & Adapted  & 0.2443 & 0.7622 & 0.2309 & 0.0057 \\
    \multirow{2}{*}{Llama-3.2-1B-Instruct}
      & Base     & 0.5009 & 0.0000 & 0.3011 & 0.0034 \\
      & Adapted  & 0.3148 & 0.8297 & 0.2377 & 0.0115 \\
    \multirow{2}{*}{Qwen3-1.7B}
      & Base     & 0.6617 & 0.0000 & 0.3317 & 0.0241 \\
      & Adapted  & 0.3441 & 0.7853 & 0.2460 & 0.0115 \\
    \bottomrule
  \end{tabular}
\end{table}

\subsection{\label{apx:primary}The primary decomposition}
Fitting \eqref{eq:model} over the $288$ runs of the three pre-registered mechanisms gives a large and precisely estimated semantic coefficient (Table~\ref{tab:mainfit}).
With distortion measured on the training text,
$\hat c = 0.647$ ($\mathrm{SE}\ 0.076$, $p < 10^{-16}$);
measured on generations,
$\hat c = 0.543$ ($\mathrm{SE}\ 0.055$, $p < 10^{-22}$).
The sign is positive throughout:
more semantic distortion,
more degradation.
The entity-targeted control is held out of this fit and analyzed separately in \apx{targeted},
since it was added after the protocol was registered.

\paragraph{Primary claim.}
The pre-registered threshold is $|c| \ge 0.10$ in standardized units.
The estimate exceeds it by a factor of six,
and the lower end of the $95\%$ interval sits at $0.498$ on the text side and $0.435$ on the generation side,
four times the threshold in both cases.
The claim is supported under both measurements of $S$.

\paragraph{The claim under the registered standard errors.}
The interval above comes from the random-intercept fit.
The registered alternative clusters on the nine model $\times$ mechanism cells,
which roughly doubles the standard error on $c$,
from $0.076$ to $0.147$ text-side and from $0.055$ to $0.120$ generation-side (Table~\ref{tab:mainfit}).
A $t$ distribution on eight degrees of freedom reads that as $p = 0.002$;
the reference suited to nine clusters carrying seven parameters is the wild cluster bootstrap with the null imposed~\citep{cameron2015practitioner},
and it gives $p = 0.060$ text-side and $p = 0.048$ generation-side,
on a $95\%$ interval that reaches the negligibility floor.
That is the resolution nine clusters supply.
The estimate itself holds at $+0.647$,
and the same dependence taken where $S$ is assigned -- the $26$ dataset $\times$ variant combinations,
within each of which $S$ is a constant -- gives standard errors of $0.108$ and $0.073$,
with bootstrap $p$ of $0.003$ and below $10^{-4}$;
the finer model $\times$ mechanism $\times$ level cells give $0.095$ and $0.078$,
with $p$ below $10^{-4}$ under both measurements.
All three are reported,
and the registered one is read as the bound its cluster count sets.

\paragraph{Dominance.}
The level slopes are $b = -0.128$,
$+0.330$ and $-0.023$ for random replacement,
sanitization and private optimization with text-side $S$,
and $-0.207$,
$+0.325$ and $-0.005$ generation-side.
The largest in magnitude is $0.330$,
against $|c| = 0.647$,
so $|c| > |b_m|$ for every mechanism under both measurements and the secondary dominance claim is supported.

\paragraph{Collinearity under the finer grid.}
The pre-registered fallback triggers at $\mathrm{VIF} > 10$ for $S$. Observed values are $8.07$ text-side and $4.43$ generation-side:
the fallback is not needed,
but the text-side figure is close enough to the threshold to deserve comment,
and it is substantially higher than the four-level design would have given.
Refitting on the original four sanitization levels alone puts the text-side VIF near $1.4$.
The finer grid is responsible.
With four widely spaced levels,
sanitization's $S$ is a violently non-monotone function of the level rank and little of it is predictable from a straight line;
with eight levels the same sigmoid is sampled densely enough that a linear term in $d$ tracks it well,
leaving less residual variation to identify $c$. Adding resolution to the mechanism's response curve made the decomposition harder to identify,
not easier -- a design lesson that runs against intuition.

\paragraph{Identifying $c$ without the level term.}
The inflation comes from the level index $d$ and $S$ being collinear,
so the direct remedy is to drop $d$ and identify $c$ from measured distortion alone.
Refitting $Y \sim \text{mechanism} + S$ without the level term or its interactions takes the text-side VIF from $8.07$ to $2.08$ and gives $\hat c = +0.403$ ($\mathrm{SE}\ 0.041$, $p < 10^{-21}$).
The estimate is smaller than the $0.647$ obtained while conditioning on $d$ -- the two specifications target different quantities -- but it is far more sharply identified and still four times the pre-registered threshold.
Whether or not the ordinal level is partialled out,
distortion carries a large,
highly significant positive coefficient;
the effect is a property of the data rather than of the collinearity between $S$ and the level.

\paragraph{The sanitization level slope inverts.}
The level slope for sanitization is positive,
meaning that after conditioning on $S$ a looser budget is associated with \textit{more} degradation.
This is a suppression artifact rather than a finding:
sanitization's $S$ is sharply non-monotone in the level rank (\S\ref{sec:cliff}),
so once $S$ absorbs the mechanism's actual damage the residual level slope inverts.
The decomposition should not be read as clean for this mechanism.

\paragraph{What identifies $c$.}
Because $S$ is constant within a condition,
the per-mechanism intercept and slope already span every linear function of $d$;
only within-mechanism curvature can identify $c$. Table~\ref{tab:residual} reports the residual of $S$ after absorbing that intercept and slope.
Between $10\%$ and $20\%$ of each mechanism's range in $S$ survives as identifying variation.
Private optimization contributes nothing on the text side,
where its $S$ is identically zero,
but contributes the largest residual share of any mechanism on the generation side.

\begin{table}[t]
  \centering
  \caption{Residual of $S$ after a per-mechanism intercept and slope in $d$. RMS is the root mean square of the residual.}
  \label{tab:residual}
  \begin{tabular}{llcccc}
    \toprule
    $S$ & Mechanism & $n$ & RMS & Range & Share \\
    \midrule
    \multirow{3}{*}{Text-side}
      & Random replacement   & 72 & 0.0292 & 0.2843 & $10.3\%$ \\
      & Sanitization         & 144 & 0.1046 & 0.5576 & $18.8\%$ \\
      & Private optimization & 72 & 0.0000 & 0.0000 & -- \\
    \midrule
    \multirow{3}{*}{Generation-side}
      & Random replacement   & 72 & 0.0300 & 0.1941 & $15.5\%$ \\
      & Sanitization         & 144 & 0.0734 & 0.4060 & $18.1\%$ \\
      & Private optimization & 72 & 0.0146 & 0.0712 & $20.4\%$ \\
    \bottomrule
  \end{tabular}
\end{table}

\subsection{\label{apx:dominated}The privacy budget is inert}
The plateau has a practical corollary,
and it was pre-registered before any outcome of the finer grid existed.
If two settings sit on the same distortion plateau and are statistically equivalent in outcome,
the looser one is dominated and can be set aside:
a stricter budget gives an at-least-as-strong guarantee at no cost.

The registered selection rule is deterministic and reads only text-side $S$,
which carries no outcome information.
Per dataset,
with $P$ the distortion at $\varepsilon = 1$,
a setting is on-plateau when $|S - P| \le 0.03P$;
STRICT is the smallest such $\varepsilon$,
PLATEAU-LOOSE the largest,
and OFF the smallest $\varepsilon$ with $S \le 0.5P$.
Applied to the eight-level grid the rule selects $\varepsilon = 0.5$,
$3$ and $6$ on \textit{both} datasets,
the plateau resolving to $\{0.5, 1, 2, 3\}$ in each case.
Tests are paired over the six dataset $\times$ model cells of seed means at the registered margin $\Delta = 0.05$.

\begin{table}[t]
  \centering
  \caption{Pre-registered equivalence tests.
  Prediction 1 expects equivalence between STRICT and PLATEAU-LOOSE;
  prediction 2 expects non-equivalence between PLATEAU-LOOSE and OFF.
  Refusal is exploratory.}
  \label{tab:tost}
  \begin{tabular}{llccc}
    \toprule
    Contrast & Metric & Mean diff. & Paired $p$ & TOST $p$ \\
    \midrule
    \multirow{4}{*}{$\varepsilon = 0.5$ vs $3$}
      & Fidelity    & $+0.0260$ & 0.106 & $0.064$ \\
      & Instruction & $+0.0048$ & 0.369 & $0.0001$ \\
      & Margin      & $+0.0102$ & 0.383 & $0.007$ \\
      & Refusal     & $+0.0021$ & 0.337 & $<0.0001$ \\
    \midrule
    \multirow{4}{*}{$\varepsilon = 3$ vs $6$}
      & Fidelity    & $-0.2115$ & 0.0005 & $0.999$ \\
      & Instruction & $-0.0927$ & 0.054 & $0.851$ \\
      & Margin      & $+0.2216$ & 0.0005 & $0.999$ \\
      & Refusal     & $-0.2669$ & 0.050 & $0.955$ \\
    \bottomrule
  \end{tabular}
\end{table}

\paragraph{Prediction 1: partially supported.}
Instruction following and the truthfulness margin pass the equivalence test (Table~\ref{tab:tost}; $p = 0.0001$ and $0.007$),
as does refusal.
Entity fidelity does not ($p = 0.064$),
missing the threshold on a mean difference of $0.026$ -- a sixth of the seed spread of that probe on the entity-poor dataset.
The registered prediction required all three,
so it holds on two of the three metrics.
Per dataset the fidelity test gives $p = 0.18$ and $p = 0.14$,
so the metric behaves the same way on both datasets,
and we report the prediction as holding on two of three metrics as registered.

\paragraph{Prediction 2: supported.}
Every metric separates PLATEAU-LOOSE from OFF,
with TOST $p$ above $0.85$ throughout,
and fidelity and the margin show significant paired differences ($p = 0.0005$).
The plateau is real and bounded:
it is not the case that all sanitization settings are interchangeable.

\paragraph{What this buys.}
A sixfold range of privacy budget is eliminated from consideration on two of three utility metrics without training anything,
since text-side distortion is computed from text alone.
Two qualifications.
The plateau lies at the high-damage end of the grid -- fidelity $0.70$ against $0.95$ past the transition -- so elimination narrows a candidate set rather than resolving the privacy--utility trade.
And the margin $\Delta = 0.05$ was chosen after inspecting the seed spread,
which the protocol records as a deviation.

\paragraph{\label{apx:clipping}Flat in the budget and in the clip.}
Training telemetry explains the flatness (Table~\ref{tab:clipping}).
Between a half and three quarters of examples hit the clipping bound at every budget,
with mean per-example gradient norms of $6.7$--$9.0$ against a bound of $1.0$.
The signal-to-noise ratio varies by a factor of $2.5$ across budgets while the clip rate does not move monotonically at all.
Since outcomes are equivalent across the same range for two of three metrics (Table~\ref{tab:tostm3}),
the binding constraint is the clipping bias,
which does not vary with the budget,
rather than the injected noise,
which does.

\begin{table}[t]
  \centering
  \caption{Private-optimization training telemetry, averaged over datasets,
  models and seeds.
  The clipping bound is $1.0$ throughout.}
  \label{tab:clipping}
  \begin{tabular}{lcccc}
    \toprule
    $\varepsilon$ & 1 & 4 & 8 & 16 \\
    \midrule
    Gradient norm    & 8.975 & 7.373 & 8.377 & 6.671 \\
    Clipped fraction & 0.753 & 0.728 & 0.561 & 0.718 \\
    Signal-to-noise  & 0.0083 & 0.0130 & 0.0164 & 0.0207 \\
    \bottomrule
  \end{tabular}
\end{table}

\begin{table}[t]
  \centering
  \caption{Equivalence across the full private-optimization budget range,
  $\varepsilon = 1$ against $16$.}
  \label{tab:tostm3}
  \begin{tabular}{lccc}
    \toprule
    Metric & Mean diff. & Paired $p$ & TOST $p$ \\
    \midrule
    Fidelity    & $-0.0377$ & 0.217 & $0.333$ \\
    Instruction & $+0.0006$ & 0.947 & $0.001$ \\
    Margin      & $+0.0167$ & 0.223 & $0.020$ \\
    \bottomrule
  \end{tabular}
\end{table}

The clipping explanation makes a prediction that can be checked directly,
and \apx{sensitivity} checks it:
varying the clipping bound over a thirtyfold range changes outcomes as little as varying the budget does.
Both knobs of the mechanism are inert at this scale;
what is not inert is the adapter rank.

\subsection{\label{apx:cliffdetail}Sanitization's response is a threshold}
The clearest evidence in the study needs no model.
Sanitization's map from privacy budget to semantic distortion is flat over half its range and then collapses (Table~\ref{tab:cliff}).
Across $\varepsilon \in \{0.5, 1, 2, 3\}$ the text-side $S$ moves by $0.0026$ on the entity-rich dataset and $0.0022$ on the entity-poor one -- a sixfold change in nominal budget leaving measured distortion where it was,
because the mechanism already resamples most tokens at each of those settings.
The collapse then happens between $\varepsilon = 4$ and $\varepsilon = 6$,
where $S$ falls by $0.349$ and $0.431$ respectively,
and by $\varepsilon = 12$ the mechanism is numerically indistinguishable from the unperturbed text.

\begin{table}[t]
  \centering
  \caption{Sanitization across the full eight-level grid.
  Text-side $S$ is flat to three decimals over the four strictest settings and collapses between $\varepsilon = 4$ and $\varepsilon = 6$. Refusal is the harmful-instruction rate.}
  \label{tab:cliff}
  \small
  \setlength{\tabcolsep}{4pt}
  \begin{tabular}{lcccccccc}
    \toprule
    $\varepsilon$ & 0.5 & 1 & 2 & 3 & 4 & 5 & 6 & 12 \\
    \midrule
    \multicolumn{9}{l}{\textbf{Entity-rich}} \\
    $S$ (text)            & 0.4597 & 0.4598 & 0.4592 & 0.4572 & 0.4322 & 0.2705 & 0.0833 & $\sim 0$ \\
    Fidelity              & 0.7195 & 0.7053 & 0.6984 & 0.7005 & 0.7210 & 0.9106 & 0.9505 & 0.9752 \\
    Instruction           & 0.1834 & 0.1826 & 0.1828 & 0.1820 & 0.1738 & 0.1873 & 0.2294 & 0.2736 \\
    Refusal               & 0.0096 & 0.0036 & 0.0047 & 0.0053 & 0.0880 & 0.2746 & 0.2756 & 0.2100 \\
    \midrule
    \multicolumn{9}{l}{\textbf{Entity-poor}} \\
    $S$ (text)            & 0.5575 & 0.5575 & 0.5576 & 0.5554 & 0.5324 & 0.3590 & 0.1012 & $\sim 0$ \\
    Fidelity              & 0.3702 & 0.3357 & 0.3551 & 0.3373 & 0.3545 & 0.3415 & 0.5103 & 0.5438 \\
    Instruction           & 0.1887 & 0.1887 & 0.1885 & 0.1805 & 0.1904 & 0.2130 & 0.3185 & 0.3307 \\
    Refusal               & 0.0000 & 0.0000 & 0.0000 & 0.0000 & 0.0000 & 0.0876 & 0.2635 & 0.3068 \\
    \bottomrule
  \end{tabular}
\end{table}

Outcomes track distortion,
not the nominal budget.
Were degradation driven by the privacy parameter,
the four leftmost columns would differ;
instead fidelity varies by $0.02$ across them on the entity-rich dataset while the budget varies sixfold,
and then moves by $0.25$ over the two columns where $S$ falls.
The entity-poor dataset shows the same shape displaced one column to the right,
consistent with its higher $S$ throughout.

\paragraph{\label{apx:plateau}The plateau carries no identifying information.}
Because the four strictest settings share a single value of $S$,
they can be used to ask what happens to alignment when distortion is held fixed and only the nominal budget moves.
Regressing $Y$ on $d$ within sanitization at $\varepsilon \le 4$,
with a per-model random intercept,
gives a slope of $-0.046$ ($\mathrm{SE}\ 0.069$, $p = 0.50$) on the entity-rich dataset and $+0.294$ ($\mathrm{SE}\ 0.058$, $p < 10^{-3}$) on the entity-poor one.
Over the same rows the seed-level standard deviation of $Y$ is $0.33$.

The two datasets therefore disagree,
and only one of them shows anything.
On the entity-rich dataset the plateau is flat in outcome as well as in distortion,
which is what a pure mediation account predicts.
On the entity-poor dataset degradation increases across a range where the measured distortion does not move,
which that account does not predict and which is evidence of a residual non-semantic component inside sanitization itself.
We report the disagreement as measured:
the effect appears on the dataset whose fidelity probe is the less precise of the two (\apx{sanity}),
and the trend is carried by instruction following,
which is measured precisely on both.

The methodological consequence is unambiguous.
Over the full sanitization range $S$ predicts $Y$ strongly on both datasets ($+1.20$, $p < 10^{-10}$ entity-rich; $+1.48$, $p < 10^{-13}$ entity-poor),
but the threshold supplies that association.
The four plateau levels cost $72$ training runs and $144$ evaluations and contribute negligible variation in the regressor of interest.
A grid designed to resolve a threshold should place its levels where the response moves,
which for this mechanism means $\varepsilon \in [4, 12]$ and not below.

\paragraph{Locating the transition without training.}
Where to place the levels can be decided before any model is trained,
from the distortion measurements alone.
Text-side $S$ holds constant across $\varepsilon \in [0.5, 4]$ -- a spread of $0.03$ on both datasets -- and then collapses over the next three settings:
the steps at $\varepsilon = 4 \to 5$ and $5 \to 6$ carry the entire change ($\Delta S = -0.16$ then $-0.19$ on the entity-rich dataset, $-0.17$ then $-0.26$ on the entity-poor one),
and by $\varepsilon = 12$ the sanitized text is at $S = 0$. The transition lives in $\varepsilon \in [4, 6]$.
A replication can read this off the distortion curve for the cost of embedding the sanitized text,
with no training at all,
and spend its budget on that interval instead of on the plateau below it.

\paragraph{Response across levels.}
Table~\ref{tab:levels} gives the remaining mechanisms against the privacy-level rank.
Random replacement is monotone in the rank on every metric,
as its linear distortion map predicts.
Private optimization stays flat throughout,
on every metric.

\begin{table}[t]
  \centering
  \caption{Metric means by mechanism and privacy level,
  averaged over models and seeds.
  Sanitization is in Table~\ref{tab:cliff}.}
  \label{tab:levels}
  \small
  \setlength{\tabcolsep}{4pt}
  \begin{tabular}{lllcccc}
    \toprule
    Metric & Dataset & Mechanism & \multicolumn{4}{c}{Strictest $\to$ loosest} \\
    \midrule
    \multirow{4}{*}{Fidelity}
      & \multirow{2}{*}{Poor} & Replacement    & 0.2563 & 0.4547 & 0.4565 & 0.5242 \\
      & & Private optim. & 0.4312 & 0.4378 & 0.4460 & 0.4650 \\
      & \multirow{2}{*}{Rich} & Replacement    & 0.8511 & 0.9287 & 0.9598 & 0.9686 \\
      & & Private optim. & 0.8964 & 0.9188 & 0.9269 & 0.9381 \\
    \midrule
    \multirow{4}{*}{Instruction}
      & \multirow{2}{*}{Poor} & Replacement    & 0.1452 & 0.1994 & 0.2288 & 0.2769 \\
      & & Private optim. & 0.3124 & 0.3192 & 0.3194 & 0.3204 \\
      & \multirow{2}{*}{Rich} & Replacement    & 0.1267 & 0.1409 & 0.1676 & 0.1820 \\
      & & Private optim. & 0.2641 & 0.2510 & 0.2541 & 0.2549 \\
    \midrule
    \multirow{4}{*}{MC1}
      & \multirow{2}{*}{Poor} & Replacement    & 0.2179 & 0.2277 & 0.2282 & 0.2388 \\
      & & Private optim. & 0.2372 & 0.2364 & 0.2358 & 0.2380 \\
      & \multirow{2}{*}{Rich} & Replacement    & 0.2207 & 0.2218 & 0.2292 & 0.2312 \\
      & & Private optim. & 0.2304 & 0.2294 & 0.2293 & 0.2279 \\
    \midrule
    \multirow{4}{*}{Margin}
      & \multirow{2}{*}{Poor} & Replacement    & $-0.7734$ & $-0.7671$ & $-0.8114$ & $-0.8527$ \\
      & & Private optim. & $-0.8356$ & $-0.8409$ & $-0.8444$ & $-0.8483$ \\
      & \multirow{2}{*}{Rich} & Replacement    & $-0.7044$ & $-0.7240$ & $-0.7498$ & $-0.7882$ \\
      & & Private optim. & $-0.8595$ & $-0.8695$ & $-0.8760$ & $-0.8802$ \\
    \bottomrule
  \end{tabular}
\end{table}

\subsection{\label{apx:components}Two dissociable pathways}
Refitting \eqref{eq:model} with each component of $Y$ as the outcome,
each z-scored so the coefficients are comparable to the composite,
shows that the composite is not merely a summary but a cancellation (Table~\ref{tab:components}).
Semantic distortion damages instruction following and entity fidelity strongly and about equally,
does little to multiple-choice truthfulness,
and \textit{improves} the truthfulness margin.

\begin{table}[t]
  \centering
  \caption{$\hat c$ per outcome component,
  both measurements of $S$. Outcomes are z-scored drops from the non-private baseline,
  so all entries are on the scale of the composite.}
  \label{tab:components}
  \begin{tabular}{lcccc}
    \toprule
    & \multicolumn{2}{c}{Text-side} & \multicolumn{2}{c}{Generation-side} \\
    \cmidrule(lr){2-3}\cmidrule(lr){4-5}
    Outcome & $\hat c$ & SE & $\hat c$ & SE \\
    \midrule
    $Y$ (composite)          & $+0.647$ & 0.076 & $+0.543$ & 0.055 \\
    Instruction following    & $+0.896$ & 0.093 & $+0.660$ & 0.070 \\
    Entity fidelity          & $+0.885$ & 0.127 & $+0.706$ & 0.094 \\
    TruthfulQA MC1           & $+0.205$ & 0.128 & $+0.288$ & 0.095 \\
    Truthfulness margin      & $-0.479$ & 0.083 & $-0.197$ & 0.065 \\
    \bottomrule
  \end{tabular}
\end{table}

Two conditions bound the interpretation.
MC1 sits at chance for every adapted model ($0.22$--$0.27$ against a $0.23$ baseline),
so its small coefficient reflects the limits of that measurement as much as the size of the effect;
it is nominally significant generation-side ($p = 0.002$) and not text-side ($p = 0.11$).
And the margin rising under distortion may reflect degraded calibration rather than improved truthfulness:
the models' errors become less confident,
which is not the same as becoming right.
The sign is robust;
its interpretation is not settled.

The practical consequence is that the composite understates the semantic pathway on the capabilities that respond,
and that any analysis run on $Y$ alone -- including the moderation test below -- inherits a cancellation that $Y$ alone conceals.

\paragraph{\label{apx:dissociation}The two routes.}
Table~\ref{tab:delta} gives each mechanism's mean change from the non-private baseline.
Private optimization costs entity fidelity as much as random replacement does on the entity-poor dataset ($0.165$ against $0.187$) while its cost to instruction following is $0.020$ against $0.126$.
Its text-side distortion is exactly zero:
it leaves the text intact.

\begin{table}[t]
  \centering
  \caption{Degradation relative to the non-private baseline,
  averaged over levels,
  models and seeds.
  Positive means a larger loss in fidelity and instruction following;
  the margin column is the rise in margin,
  i.e.\ less confident error.}
  \label{tab:delta}
  \begin{tabular}{llccc}
    \toprule
    Dataset & Mechanism & Fidelity & Instruction & Margin \\
    \midrule
    \multirow{3}{*}{Poor}
      & Replacement    & $0.1873$ & $0.1257$ & $+0.1094$ \\
      & Sanitization   & $0.2167$ & $0.1134$ & $+0.2461$ \\
      & Private optim. & $0.1652$ & $0.0204$ & $+0.0682$ \\
    \midrule
    \multirow{4}{*}{Rich}
      & Replacement    & $0.0450$ & $0.1090$ & $+0.2128$ \\
      & Sanitization   & $0.1744$ & $0.0640$ & $+0.2831$ \\
      & Entity-targeted& $0.0044$ & $0.0023$ & $+0.1367$ \\
      & Private optim. & $0.0520$ & $0.0073$ & $+0.0831$ \\
    \bottomrule
  \end{tabular}
\end{table}

The mechanisms therefore damage different capabilities:
the text-perturbing family costs instruction following and moves the truthfulness margin,
while private optimization costs factual fidelity without semantic distortion of the training signal.
A single one-dimensional degradation axis does not describe this.

\paragraph{A floor for generation-side distortion.}
Private optimization's generation-side $S$ is not zero,
which is what keeps a mediation account available for it.
Interpreting that number requires knowing how much generation-side distortion arises from nothing at all.
The six-epoch runs of \apx{sensitivity} provide it:
they are trained on the same clean text as the non-private baseline and differ only in epoch count,
so their generation-side score of $0.127$ is pure configuration and seed noise.
Against that floor,
entity-targeted sanitization sits at $0.128$ -- indistinguishable -- private optimization clears it by $0.03$--$0.04$ at $0.164$,
random replacement reaches $0.184$,
and sanitization on its plateau reaches $0.391$,
three times the floor.
Private optimization's generation-side signal is real but small,
and it is the cleanest available estimate of the non-semantic route.

\paragraph{Fabrication mediates the semantic pathway.}
The semantic route is not only real but decomposable:
much of what distortion costs,
it costs by making the model fabricate.
We fit a mediation over the registered mechanisms~\citep{baron1986moderator, imai2010general} with the causal ordering distortion $\to$ fabrication rate $\to$ degradation,
taking the outcome to be the mean of the instruction-following and truthfulness $z$-scores.
Entity fidelity is deliberately excluded from this outcome:
its correctness component is mechanically tied to the fabrication mediator,
and leaving it in would manufacture the mediation we are testing for.
In standardized units,
a unit of text-side distortion raises the fabrication rate by $1.14$;
conditioning on fabrication,
the direct effect of distortion on the outcome falls from $+0.518$ to $+0.269$.
The indirect effect is $+0.250$ (percentile bootstrap $95\%$ CI $[+0.173, +0.339]$, $2000$ resamples~\citep{preacher2008asymptotic}),
so fabrication carries about $48\%$ of the total.
The mediated share is concentrated on the entity-poor dataset,
where distortion drives fabrication hardest;
on the entity-rich dataset the outcome barely moves with distortion once entity fidelity is set aside,
leaving little for any channel to carry.
The decomposition of \apx{primary} establishes that distortion matters;
this establishes how it does its damage.

\paragraph{Fabrication, not corruption.}
The extraction probe separates three outcomes:
a correct answer,
a \texttt{corrupted} answer whose values all appear in the source record but do not match the gold set,
and a \texttt{fabricated} answer containing at least one value absent from the record.
The distinction matters because only the third is hallucination in the usual sense -- the model inventing content rather than selecting the wrong content.

Semantic distortion acts predominantly on the fabrication channel.
Fitting \eqref{eq:model} in raw rate units gives $\hat c = +0.255$ ($p < 10^{-14}$) for fabrication against $\hat c = +0.032$ ($p < 10^{-4}$) for corruption with text-side distortion,
and $+0.217$ against $+0.024$ generation-side.
Distortion multiplies invention roughly eightfold more than it multiplies selection error.

\begin{table}[t]
  \centering
  \caption{Fabrication rate by mechanism and privacy level,
  entity-rich dataset.
  The non-private baseline is $0.0042$.}
  \label{tab:fabrication}
  \footnotesize
  \setlength{\tabcolsep}{3pt}
  \begin{tabular}{lcccccccc}
    \toprule
    Mechanism & \multicolumn{8}{c}{Strictest $\to$ loosest} \\
    \midrule
    Random replacement   & 0.0643 & 0.0184 & 0.0100 & 0.0060 & & & & \\
    Sanitization         & 0.1685 & 0.1727 & 0.1833 & 0.2011 & 0.1821 & 0.0477 & 0.0214 & 0.0033 \\
    Entity-targeted      & 0.0042 & 0.0039 & & & & & & \\
    Private optimization & 0.0377 & 0.0214 & 0.0163 & 0.0124 & & & & \\
    \bottomrule
  \end{tabular}
\end{table}

On the entity-rich dataset the effect is stark (Table~\ref{tab:fabrication}).
The non-private baseline fabricates on $0.4\%$ of records.
Uniform sanitization fabricates on $16.9$--$20.1\%$ across its five plateau settings -- a factor of forty to forty-eight -- and then returns to $0.3\%$ once past the transition.
Random replacement and private optimization reach $6.4\%$ and $3.8\%$ at their strictest,
an order of magnitude less.
Entity-targeted sanitization fabricates at $0.42\%$ and $0.39\%$:
at the baseline rate,
despite displacing entities as hard as the uniform mechanism does.

That last row is the cleanest statement of the study's thesis in the vocabulary of hallucination.
Privacy mechanisms that damage the semantics of the text around the sensitive content cause models to invent facts;
a mechanism that damages the sensitive content just as hard while leaving its surroundings intact does not.
Invention tracks corruption of the carrier,
and the entity spans can be displaced on their own without provoking it.

One property of the scorer:
an empty response is counted as fabrication,
since it contains no value drawn from the record.
Empty responses are rare in these conditions,
but the fabrication rates above are therefore upper bounds.

\subsection{\label{apx:targeted}Entity-targeted sanitization separates the two routes}
Uniform sanitization damages entities and the text around them together,
which is why its two distortion components move together within that family.
The entity-targeted variant resamples each distinct entity once and substitutes it everywhere in both fields,
leaving the carrier text verbatim.
It achieves the same entity displacement as uniform sanitization at matched budget -- $S_{\text{entity}} = 0.669$ against $0.671$ -- with a quarter of the collateral damage to the carrier ($0.111$ against $0.502$).

The outcome difference is the largest single contrast in the study (Table~\ref{tab:targeted}).
At matched entity damage,
the targeted variant tracks the clean baseline:
fidelity $0.967$ against a clean baseline of $0.975$ and a uniform-sanitization value of $0.705$;
instruction following $0.260$ against $0.264$ and $0.183$;
fabrication $0.0042$,
identical to the clean baseline's $0.0042$,
against $0.173$.
Refusal is likewise preserved -- $0.209$ against a baseline $0.196$,
where uniform sanitization at the same budget gives $0.004$.
Both budgets behave identically,
the mechanism having its own plateau.

\paragraph{The decomposition now identifies.}
Text-side distortion splits into an entity-retention term and a carrier term measured after masking entity spans.
Across the uniform mechanisms alone the two are collinear at $r = 0.991$,
giving VIF of $58.4$ -- six times the threshold at which the protocol turns to the cross-mechanism contrast -- so the split rests on a further condition.
The targeted variant is precisely the observation that breaks the collinearity,
since it is the only condition with high entity damage and low carrier damage.
Including it drops the correlation to $r = 0.716$ and the VIF to $2.1$,
and the decomposition becomes estimable:
\begin{equation*}
  \hat c_{\text{carrier}} = +1.83,\qquad
  \hat c_{\text{entity}} = -0.72,
\end{equation*}
both with $p < 10^{-4}$ over the $198$ entity-rich text-perturbing runs.
Holding carrier damage fixed,
damaging entities harder does not degrade alignment;
the degradation is carried by what happens to the surrounding text.
This is the study's sharpest mechanistic statement,
and it inverts the intuition that motivated the entity probe in the first place.

Two qualifications.
The negative entity coefficient is estimated from a contrast between two mechanism families rather than within one,
so it inherits whatever else distinguishes them.
And the targeted variant is the only source of the identifying variation,
so the estimate rests on two conditions and a third condition would be needed to check it.

\subsection{\label{apx:refusal}Refusal behavior}
A separate suite measured refusal on the $520$ harmful instructions of AdvBench~\citep{zou2023universal} and the $450$ prompts of XSTest~\citep{rottger2024xstest},
split between benign-but-alarming and genuinely unsafe.
The second set is a control:
it separates a model that has become discriminating from one that refuses indiscriminately.
Empty-response rate is $0.000$ on the benign-unsafe set and at most $0.023$ on harmful instructions,
and mean response length stays between $54$ and $108$ words across mechanisms,
so no result below is an artifact of degenerate generation.

\paragraph{Adaptation destroys refusal before any mechanism acts.}
Table~\ref{tab:refusalbase} shows the untuned models against their non-private adapted counterparts.
Refusal on harmful instructions falls by $0.47$,
$0.49$ and $0.76$.
One model retains $2\%$ of its refusal behavior after three epochs on a benign extraction dataset.
This dwarfs every mechanism effect that follows and is the correct baseline for reading them.

\begin{table}[t]
  \centering
  \caption{Refusal before and after non-private adaptation.}
  \label{tab:refusalbase}
  \begin{tabular}{llcccc}
    \toprule
    Model & State & Harmful & Benign & Unsafe & Discrim. \\
    \midrule
    \multirow{2}{*}{Falcon3-1B-Instruct}
      & Base    & 0.9712 & 0.2360 & 0.9550 & 0.7190 \\
      & Adapted & 0.5061 & 0.0180 & 0.4858 & 0.4678 \\
    \multirow{2}{*}{Llama-3.2-1B-Instruct}
      & Base    & 0.6558 & 0.0360 & 0.3000 & 0.2640 \\
      & Adapted & 0.1612 & 0.0033 & 0.1767 & 0.1733 \\
    \multirow{2}{*}{Qwen3-1.7B}
      & Base    & 0.7827 & 0.0480 & 0.5850 & 0.5370 \\
      & Adapted & 0.0224 & 0.0000 & 0.0775 & 0.0775 \\
    \bottomrule
  \end{tabular}
\end{table}

\paragraph{Mechanisms move in opposite directions.}
Text perturbation erodes refusal further;
private optimization raises it well above the non-private baseline (Table~\ref{tab:refusalmech}).
On the entity-rich dataset the private effect is monotone in the budget -- $0.369$,
$0.333$,
$0.315$,
$0.289$ from strictest to loosest -- so a tighter guarantee preserves more.
The mechanism is plausible:
noise and clipping limit how far the adapted weights travel from the base model,
and the base model's refusal behavior travels with them.

\begin{table}[t]
  \centering
  \caption{Refusal by mechanism, averaged over levels,
  models and seeds.}
  \label{tab:refusalmech}
  \begin{tabular}{llcccc}
    \toprule
    Dataset & Mechanism & Harmful & Benign & Unsafe & Discrim. \\
    \midrule
    \multirow{4}{*}{Poor}
      & Baseline       & 0.2643 & 0.0040 & 0.2494 & 0.2454 \\
      & Replacement    & 0.1450 & 0.0018 & 0.1912 & 0.1895 \\
      & Sanitization   & 0.0822 & 0.0008 & 0.0821 & 0.0813 \\
      & Private optim. & 0.3972 & 0.0140 & 0.2686 & 0.2546 \\
    \midrule
    \multirow{5}{*}{Rich}
      & Baseline       & 0.1955 & 0.0102 & 0.2439 & 0.2337 \\
      & Replacement    & 0.1599 & 0.0293 & 0.3254 & 0.2961 \\
      & Sanitization   & 0.1089 & 0.0108 & 0.1764 & 0.1656 \\
      & Entity-targeted& 0.2007 & 0.0067 & 0.2678 & 0.2611 \\
      & Private optim. & 0.3267 & 0.0396 & 0.2853 & 0.2457 \\
    \bottomrule
  \end{tabular}
\end{table}

\paragraph{The threshold again, and it is total.}
At all five sanitization settings on the distortion plateau,
models trained on the entity-poor dataset refuse \textit{nothing}:
harmful-instruction refusal is exactly $0.0000$ and discrimination exactly $0.0000$,
across three models,
three seeds and five budgets.
On the entity-rich dataset the same cells read $0.0036$ to $0.0096$.
Past the transition,
refusal returns above the non-private baseline (Table~\ref{tab:cliff}).
Safety follows distortion,
rather than the budget,
at the same threshold that governs utility -- and the four extra levels turn what was a two-point observation into an unbroken run of five.

\paragraph{Distortion predicts refusal loss.}
Fitting \eqref{eq:model} with refusal as the outcome,
in raw rate units rather than standardized,
gives $\hat c = -0.101$ ($p < 10^{-3}$) on the text side and $-0.127$ ($p < 10^{-9}$) generation-side for harmful-instruction refusal,
and $-0.116$ and $-0.137$ for discrimination.
One standard deviation of semantic distortion costs roughly twelve points of refusal rate.
The pathway that damages instruction following and fidelity damages refusal in the same direction.

\paragraph{The control fires.}
Private optimization's gain is not purely discrimination.
On the entity-rich dataset it raises refusal on \textit{benign} prompts from $0.0102$ to $0.0396$,
and to $0.0573$ at the strictest budget,
while discrimination rises only from $0.2337$ to $0.2457$.
A substantial share of the extra refusal is indiscriminate.
Reported alone,
the harmful-instruction number would license the claim that private optimization restores safety;
with the control,
the supportable claim is that it preserves refusal behavior partly by making the model refuse more in general.

\paragraph{The two probes agree.}
Harmful-instruction refusal and discrimination correlate at $r = 0.858$ across conditions,
so the two instruments are measuring a common construct.
Refusal correlates moderately with instruction following ($r = 0.449$) and not at all with entity fidelity ($r = 0.025$),
and within mechanisms the correlation ranges from $-0.77$ to $+0.47$.
Safety and capability degrade largely separably.

\subsection{\label{apx:composition}Robustness: composition, scale, and moderation}
The main grid holds the mechanisms apart.
To ask what a deployment pays when it both sanitizes its data \textit{and} trains privately,
we lift that restriction and train with DP-SGD on the already-sanitized datasets:
two sanitization budgets ($\varepsilon = 3, 6$) crossed with two privacy budgets ($\varepsilon = 4, 8$),
on both datasets across the three models and three seeds -- $72$ composed adapters.
Each composed cost is set against the two mechanisms run alone at the same budgets,
against their sum (costs add) and against the larger of the two (one dominates).

The composed cost lands between those two null models.
Summed over the three capability probes,
the composed degradation is $+0.306$,
against $+0.372$ if the costs added exactly and $+0.258$ for whichever single mechanism costs more.
The gap to the additive prediction is $-0.066$ (percentile bootstrap $95\%$ CI $[-0.104, -0.028]$~\citep{efron1994introduction}) and the gap to the larger single cost is $+0.048$:
composing reliably costs more than either mechanism alone,
and slightly less than their sum.
The two penalties largely stack,
with a mild saturation rather than a synergy -- a deployment that both sanitizes and trains privately pays close to the sum of the two,
not more.

Where they stack follows the pathway split of \apx{dissociation}.
Instruction following is sanitization's to damage:
the composed cost ($0.066$) is indistinguishable from sanitization alone ($0.073$; observed minus larger-single $-0.007$),
and private optimization adds nothing to it,
so the budget inertness of \apx{primary} survives composition.
Entity fidelity is damaged by both and roughly additively -- composed $0.240$ against a sum of $0.278$ and a larger-single of $0.193$ -- and truthfulness moves for neither.
Composition therefore springs no surprise:
it inherits sanitization's instruction cost and adds private optimization's fidelity cost on top,
with no interaction beyond the saturation expected of bounded metrics.

\paragraph{\label{apx:scale}The pathway at larger scale.}
The main grid runs at $1$--$2$B parameters.
As a two-point check on whether the semantic pathway is an artifact of that range,
we replicate it at $8$B (Qwen3-8B):
$62$ adapters spanning both datasets,
all five mechanism families and up to three seeds,
with the sanitization and replacement grids reduced to keep the check affordable.
Fitting the primary decomposition on these runs alone gives a text-side distortion coefficient of $\hat c = +0.376$ ($\mathrm{SE}\ 0.189$, $p = 0.052$, $n = 53$),
against $+0.63$ for the $1$--$2$B fit under the same ordinary-least-squares specification.

The pathway persists:
the coefficient keeps its sign and its order of magnitude,
attenuating to roughly $60\%$ of the smaller-model value on a grid reduced to what a scale check can afford.
That is the direction a scale story predicts,
a larger model's priors absorbing more of what a small adaptation set distorts,
and with two points we read it as evidence that the pathway survives at this scale,
leaving the shape of its decay to a denser grid.
Generation-side distortion was not scored for these adapters,
so the check is text-side only.

\paragraph{\label{apx:moderation}Moderation by entity density.}
Extending \eqref{eq:model} with an $S \times$ dataset interaction tests whether distortion costs more where text is fact-dense.
On the pre-registered composite outcome the two measurements of $S$ give different answers:
$+0.047$ ($p = 0.37$, $95\%$ interval $[-0.056, +0.149]$) text-side,
and $+0.153$ ($p = 0.006$, $[+0.044, +0.262]$) generation-side against a main effect of $0.405$.
The registered primary is the text-side measurement,
on which the interaction is indistinguishable from zero.

Beneath the disagreement the components disagree more sharply,
and each moves in the direction its own measurement noise predicts (Table~\ref{tab:moderation}).
Entity fidelity is $3.4\times$ noisier on the entity-poor dataset and yields a large positive interaction;
instruction following,
the only component measured with comparable precision on both datasets,
yields a negative one under both measurements of $S$. Differential noise attenuates the estimated coefficient on whichever dataset is noisier,
so the fidelity result is what an artifact would produce,
and it is large enough to drive the composite by itself.

\begin{table}[t]
  \centering
  \caption{$S \times$ dataset interaction by component,
  outcomes z-scored.
  Positive favors the entity-rich dataset,
  as the hypothesis predicts.
  Seed spread is the mean within-condition range on each dataset.}
  \label{tab:moderation}
  \begin{tabular}{lccc}
    \toprule
    Component & Seed spread (poor / rich) & Text-side & Gen.-side \\
    \midrule
    Composite $Y$         & --             & $+0.047$ & $+0.153$ \\
    Entity fidelity       & $0.175 / 0.041$ & $+0.477$ & $+0.611$ \\
    Instruction following & $0.029 / 0.036$ & $-0.199$ & $-0.163$ \\
    Truthfulness margin   & $0.032 / 0.062$ & $-0.053$ & $+0.087$ \\
    \bottomrule
  \end{tabular}
\end{table}

The reading is that the registered measurement places moderation by entity density at zero,
that the generation-side significance is attributable to the one component whose instrument differs in precision across datasets,
and that the component measured comparably points the other way.
The methodological consequence generalizes:
testing moderation \textit{by dataset} requires an outcome measurable with equal precision \textit{on both datasets},
which the extraction probe supplies on the entity-rich dataset alone.

\subsection{\label{apx:leakage}Leakage}
The memorization probe sits below its detection floor.
Record-level rates are $0.0115$ for base models,
$0.0133$ for the non-private baselines,
and $0.0033$--$0.0047$ across all mechanisms.
Every privacy mechanism reduces the rate relative to non-private adaptation,
which is directionally sensible,
but standard deviations are as large as the means ($0.004$--$0.011$),
so no pair is distinguishable and the ordering between the base models and the non-private baselines even inverts.

The manipulation check therefore operates at its detection floor at this scale,
and the privacy side of the trade rests on the R\'enyi accounting.
The floor itself is a measurement result.
Counting every string match puts the rate at $0.15$--$0.21$;
restricting to distinctive spans -- at least eight characters,
absent from the visible half of the record -- puts it at the values above,
which locates the bulk of an unfiltered extraction signal in chance matches on dates and common names.

\subsection{Scope of the measurements}
\paragraph{What each outcome component contributes.}
Multiple-choice truthfulness sits at chance for every adapted model.
It contributes noise to the composite,
widening intervals while leaving $c$ unbiased,
and its own coefficient reaches significance under one measurement of $S$,
which is the signature of an uninformative component.

\paragraph{Where the extraction probe is exact.}
Gold span annotations make the probe exact on the entity-rich dataset,
which is where the fidelity conclusions are drawn.
On the entity-poor dataset its seed spread exceeds the effects being measured,
so that cell is read as bounded by noise:
it is the metric that carries the moderation estimate,
and the one pre-registered equivalence test falling outside its margin does so here.

\paragraph{Leakage sits at the probe's detection floor.}
Measured leakage stays at noise level everywhere (\apx{leakage}),
so the privacy side of the trade rests on the formal guarantees.

\paragraph{Distortion under a second encoder.}
All primary values of $S$ come from a single sentence encoder (\texttt{BAAI/\allowbreak bge-base-\allowbreak en-v1.5}~\citep{xiao2024cpack}).
To check that the conclusions do not depend on it we recomputed every distortion score with a second,
architecturally unrelated encoder (\texttt{all-\allowbreak MiniLM-\allowbreak L6-v2}~\citep{wang2020minilm, reimers2019sbert}) and refit the primary model against the result;
the outcome is Table~\ref{tab:encoder}.
The two encoders disagree on the \textit{scale} of distortion -- MiniLM assigns uniformly larger distances,
with $S_\text{text}$ ranging to $0.91$ where the primary encoder reaches $0.56$ -- but agree on its \textit{ordering}:
Pearson correlation $0.997$ on the text side and $0.999$ on the generation side,
Spearman $0.997$ and $0.998$.
Since $S$ enters the model z-scored,
the scale difference is absorbed and the shared coefficient barely moves:
$\hat c$ shifts from $+0.647$ to $+0.608$ text-side and from $+0.543$ to $+0.548$ generation-side,
both changes smaller than one standard error,
with every per-mechanism level slope and both pre-registered inequalities preserved.
One component of $S$ is invariant to this check by construction:
the entity score $S_\text{entity}$ is a token-membership rate rather than an embedding distance,
so it is identical under the two encoders and does not rest on either.
What a different encoder could still move is the \textit{location} of the threshold,
which the ordering agreement does not pin down and which the pre-registered selection rule of \apx{dominated} consumes directly.

\begin{table}[h]
\centering
\begin{tabular}{@{}lccc@{}}
\toprule
Fit & Primary $\hat c$ (SE) & MiniLM $\hat c$ (SE) & VIF \\
\midrule
Text-side $Y$       & $+0.647\ (0.076)$ & $+0.608\ (0.083)$ & $8.07 \to 9.08$ \\
Generation-side $Y$ & $+0.543\ (0.055)$ & $+0.548\ (0.057)$ & $4.43 \to 4.66$ \\
\bottomrule
\end{tabular}
\caption{The shared distortion coefficient under the two encoders,
$n = 288$ registered runs each.
Rank correlation of $S$ between encoders is $0.997$ (text) and $0.999$ (generation).
Both coefficients clear the pre-registered $\Delta = 0.10$ and both preserve $|c| > |b_m|$ for every mechanism.}
\label{tab:encoder}
\end{table}

\paragraph{What identifies the decomposition.}
The VIF on $S$ is $8.07$ text-side,
against a fallback threshold of $10$,
and the entity--carrier decomposition of \apx{targeted} is identified by a single mechanism family occupying a single corner of the distortion space.
A mechanism occupying that corner differently would sharpen both estimates.

\paragraph{Input--output correspondence under perturbation.}
Both uniform text mechanisms rewrite the instruction field and the target field independently,
with separate random keys.
For the extraction template this means the gold answer is resampled differently from its occurrence in the document,
so the training target is no longer findable in the input.
In clean text $95.8\%$ of extraction examples contain their answer verbatim in the document;
under sanitization at $\varepsilon = 1$ this falls to $40.6\%$ and under replacement at $p = 0.50$ to $25.4\%$.
Part of the fidelity damage attributed to semantic distortion is therefore attributable to the task becoming partly unlearnable,
which is a distinct causal route.
This is faithful to what a naive field-by-field pipeline would do in practice,
and the correlational claims of \apx{primary} are unaffected -- $S$ still predicts degradation -- but the mechanism of that degradation is not purely semantic drift.
The entity-targeted variant preserves correspondence by construction,
which is a second reason its contrast in \apx{targeted} is informative,
and it bounds that contrast to carrier damage together with the correspondence it retains.

\paragraph{Scale.}
Three base models between $1$B and $1.7$B parameters,
LoRA adapters,
three epochs.
The rank sensitivity of \apx{sensitivity} shows that adapter capacity alone moves outcomes more than either private-optimization knob,
so the clipping saturation in particular should be read as specific to this scale and configuration.

\paragraph{Exploratory analyses.}
The component decomposition (\apx{components}),
the entity-targeted control (\apx{targeted}),
the refusal suite (\apx{refusal}),
the configuration blocks (\apx{sensitivity}) and the telemetry analysis (\apx{clipping}) were not pre-registered.
They are reported as exploratory;
the confirmatory claims are those of \apx{primary},
\apx{moderation} and \apx{dominated}.

\subsection{\label{apx:deviations}Deviations from the pre-registered analysis}
\paragraph{Reference for generation-side distortion.}
The protocol specifies distortion between generations of the adapted and the \textit{base} model.
The implementation measures against the non-private adapted model of the same base model,
dataset and seed.
The latter isolates the mechanism's contribution from the cost of adaptation itself,
and is consistent with $Y$ being defined against the same reference,
but it is not what was registered.

\paragraph{Standard errors.}
The protocol specifies cluster-robust standard errors~\citep{cameron2015practitioner} by model $\times$ mechanism cell.
Reported fits use a per-model random intercept,
which carries the same dependence in a variance component.
The registered errors are reported alongside them in Table~\ref{tab:mainfit} and discussed in \apx{primary}.

\paragraph{Equivalence margin.}
The margin $\Delta = 0.05$ used in \apx{dominated} was selected after inspecting the seed spread rather than in advance,
and is recorded as post-hoc in the pre-registration itself.

\paragraph{Level count.}
The protocol describes four settings per mechanism.
Sanitization was extended to eight after the first pass showed its response was unresolved,
and the entity-targeted control contributes two.
The rank coding is unchanged -- levels are ranked within a mechanism and z-scored within it,
exactly as registered -- but the sanitization ranks now run $1$ to $8$,
and $b_{\text{sanitization}}$ is consequently a slope per standard deviation of an eight-point rank rather than a four-point one.

\paragraph{Scale for the entity-targeted control.}
The control postdates the registration,
so it is excluded from the confirmatory fit of \apx{primary}.
Its $Y$ is computed on the scale fixed by the three registered mechanisms -- their mean and standard deviation -- rather than by recomputing the standardization over an enlarged pool,
so that every registered value is numerically unchanged by the control's presence.

\subsection{\label{apx:mainfits}Primary and entity-targeted fits}
The two tables the main text refers to are reproduced here in full.

\begin{table}[t]
  \centering
  \caption{Estimates of \eqref{eq:model}, per-mechanism intercept $a_m$ and per-model random intercept.
  $d$ is the privacy-level rank within a mechanism,
  z-scored;
  higher $d$ is a looser setting.
  Est.,
  SE and $p$ are the random-intercept fit.
  CR~SE is the registered cluster-robust standard error over the nine model $\times$ mechanism cells,
  and $p_{\text{wcb}}$ its wild cluster bootstrap $p$-value,
  Webb weights and $10^5$ replicates with the null imposed,
  applied to the slope hypotheses.
  Both come from the ordinary least squares fit of the same specification,
  whose estimates coincide with the tabulated ones text-side and differ by at most $0.02$ generation-side,
  where $\hat c = +0.5239$.}
  \label{tab:mainfit}
  \begin{tabular}{lccccc}
    \toprule
    Term & Est. & SE & $p$ & CR SE & $p_{\text{wcb}}$ \\
    \midrule
    \multicolumn{6}{l}{\textit{Text-side} $S$, VIF on $S$ $= 8.07$} \\
    $a$, random replacement     & $+0.5564$ & 0.1183 & $<10^{-5}$ & 0.0538 & -- \\
    $a$, sanitization           & $-0.3806$ & 0.1230 & 0.002 & 0.1037 & -- \\
    $a$, private optimization   & $+0.5023$ & 0.1401 & $<10^{-3}$ & 0.3664 & -- \\
    $b$, random replacement     & $-0.1277$ & 0.0622 & 0.040 & 0.0677 & 0.142 \\
    $b$, sanitization           & $+0.3300$ & 0.0718 & $<10^{-4}$ & 0.1504 & 0.035 \\
    $b$, private optimization   & $-0.0225$ & 0.0410 & 0.583 & 0.0402 & 0.886 \\
    $c$, semantic distortion    & $+0.6473$ & 0.0762 & $<10^{-4}$ & 0.1473 & 0.060 \\
    Model variance              & $0.0316$ & -- & -- & -- & -- \\
    \midrule
    \multicolumn{6}{l}{\textit{Generation-side} $S$, VIF on $S$ $= 4.43$} \\
    $a$, random replacement     & $+0.6002$ & 0.1226 & $<10^{-6}$ & 0.0779 & -- \\
    $a$, sanitization           & $-0.2412$ & 0.1192 & 0.043 & 0.0592 & -- \\
    $a$, private optimization   & $+0.1846$ & 0.1253 & 0.141 & 0.3249 & -- \\
    $b$, random replacement     & $-0.2070$ & 0.0554 & $<10^{-3}$ & 0.0726 & 0.081 \\
    $b$, sanitization           & $+0.3254$ & 0.0639 & $<10^{-4}$ & 0.1274 & 0.032 \\
    $b$, private optimization   & $-0.0050$ & 0.0397 & 0.901 & 0.0354 & 0.891 \\
    $c$, semantic distortion    & $+0.5431$ & 0.0553 & $<10^{-4}$ & 0.1203 & 0.048 \\
    Model variance              & $0.0351$ & -- & -- & -- & -- \\
    \bottomrule
  \end{tabular}
\end{table}

\begin{table}[t]
  \centering
  \caption{Entity-targeted against uniform sanitization at matched budget,
  and the non-private baseline,
  on the entity-rich dataset.
  Entity displacement is matched;
  everything else is not.}
  \label{tab:targeted}
  \begin{tabular}{lccc}
    \toprule
    & Baseline & Uniform & Targeted \\
    \midrule
    $S$ (text)            & 0.0000 & 0.4598 & 0.1107 \\
    $S_{\text{entity}}$   & 0.0000 & 0.6708 & 0.6685 \\
    $S_{\text{carrier}}$  & 0.0000 & 0.5016 & 0.1117 \\
    $S$ (generation)      & --    & 0.3471 & 0.1296 \\
    \midrule
    Fidelity              & 0.9746 & 0.7053 & 0.9665 \\
    Instruction           & 0.2639 & 0.1826 & 0.2598 \\
    Fabrication           & 0.0042 & 0.1727 & 0.0042 \\
    Refusal               & 0.1955 & 0.0036 & 0.2088 \\
    \bottomrule
  \end{tabular}
\end{table}

The level slopes are $-0.128$,
$+0.330$ and $-0.023$ for replacement,
sanitization and private optimization text-side,
and $-0.207$,
$+0.325$ and $-0.005$ generation-side.
The largest is $0.330$ against $|c| = 0.647$,
so $|c| > |b_m|$ holds for every mechanism under both measurements and the dominance claim is supported.
The private optimizer's own slope is the sharpest case:
$-0.023$ ($\mathrm{SE}\ 0.041$),
so tightening $\varepsilon$ from $16$ to $1$ moves degradation by nothing we can detect.
Private training costs alignment,
but as a fixed penalty arriving with the mechanism rather than a gradient to trade against.
Identifying $c$ from distortion alone,
dropping the collinear level term,
takes the VIF from $8.07$ to $2.08$ and gives $\hat c = +0.403$ ($\mathrm{SE}\ 0.041$, $p < 10^{-21}$):
the two specifications target different quantities,
but distortion carries a large,
sharply identified coefficient either way.

\begin{figure}[t]
  \centering
  \includegraphics[width=0.78\linewidth]{threshold.pdf}
  \caption{\textbf{Half of the registered sanitization levels are behaviorally inert.} \textbf{Left:} distortion of the sanitized text against the budget.
  \textbf{Right:} the resulting degradation,
  error bars over three seeds.
  The shaded band is the plateau,
  $\varepsilon \le 3$,
  where distortion varies by $0.003$;
  the transition falls in $\varepsilon \in [4,
  6]$ and can be located from the distortion curve alone,
  before training anything.}
  \label{fig:threshold}
\end{figure}
